\section{Progress and Implemented Work}

The team has divided the current implementation into data structure, frontend, and image algorithm work. Each part is designed to be independently developed and then integrated through shared data fields and documented interfaces.

\subsection{Data Structure Module Progress}

% Data structure teammate: describe completed SQL schema, simulation data
% generator, CSV generation, import workflow, interface documentation, and evidence.

\subsection{Frontend Module Progress}

% Frontend teammate: describe implemented pages/components, UI state,
% mocked or real data connection, screenshots, and remaining frontend work.

\subsection{Image Algorithm Module Progress}

The image algorithm module has been implemented as a CPU-based YOLOv8n baseline under the repository folder \texttt{image\_algorithm/} or the corresponding image algorithm folder in the shared project. Its purpose is to convert local images, or frames extracted from a future video stream, into structured seat status records.

The implemented module includes:

\begin{itemize}
  \item \texttt{main.py} as the command-line entry point;
  \item \texttt{src/detector.py} for YOLO model loading and inference;
  \item \texttt{src/geometry.py} for detection box and seat ROI matching;
  \item \texttt{src/state.py} for converting person/object signals into seat states;
  \item \texttt{src/storage.py} for Redis and MySQL integration;
  \item \texttt{src/pipeline.py} for the complete image-to-status workflow;
  \item \texttt{config.yaml}, \texttt{school\_test\_config.yaml}, and \texttt{web\_samples\_config.yaml} for reproducible configuration;
  \item \texttt{setup.ps1}, \texttt{requirements.txt}, and \texttt{README.md} for teammate setup;
  \item unit tests under \texttt{tests/}.
\end{itemize}

The module supports two modes. In normal mode, it reads \texttt{seat\_config} from MySQL, writes the latest seat states to Redis, and inserts historical rows into MySQL \texttt{seat\_status\_log}. In local demonstration mode, it can run with \texttt{--dry-run --skip-mysql --skip-redis}, use fallback seats from YAML, and generate \texttt{outputs/latest\_seat\_status.csv}. This allows the image logic to be tested even when local database services are not configured.

The current state rules are:

\begin{itemize}
  \item inactive seat: \texttt{unavailable};
  \item person detected in ROI: \texttt{occupied}, with \texttt{suspect\_duration = 0};
  \item no person but object detected: accumulate \texttt{suspect\_duration}; if it reaches the threshold, output \texttt{suspected};
  \item no person and no object: \texttt{free}, with \texttt{suspect\_duration = 0}.
\end{itemize}

The module was uploaded to the shared GitHub repository in commit \texttt{3a99c09} with the message \texttt{Add image algorithm module}. The upload intentionally excluded virtual environments, model weights, local password files, real school images, third-party sample images, and generated output files.
